\documentclass[11pt]{article}
\usepackage{fontspec}
\usepackage{booktabs}
\usepackage{hyperref}
\setmainfont{Libertinus Sans}
\title{Old Chinese Tibetan Transcription Scheme}
\author{Kelvin Yang}
\date{August 2026}
\begin{document}
\maketitle

This is a scheme for spelling Old Chinese in Tibetan script, based on the Baxter--Sagart reconstruction and intended to be concise and natural. It is primarily recreational.

\section{Initials}
Ordinary initials borrow Tibetan letters of the nearest value or series. Uvulars use the Tibetan palatal series. Abstract $C$, $N$ and $S$ are reserved for consonants whose value remains unspecified.

\section{Minor syllables and preinitials}
Material before the Baxter--Sagart dot forms a short preceding stack. An explicit vowel takes its ordinary sign; absence of a written vowel remains distinct from explicit $ə$. Tightly attached preinitials join the main initial in one stack.

\section{Pharyngealization, r and labialization}
Pharyngealization $ˤ$, post-initial $r$ and labialization $ʷ$ use a fixed set of subjoined signs. $P$-$rV$ and $PrV$ share one spelling, while initial $rV$ retains a base r.

\section{Vowels and codas}
Lowercase $a$ is unmarked and capital $A$ is written explicitly. Other vowels and all codas use the fixed correspondences given in the web edition. Stop codas follow ordinary Tibetan suffix practice.

\section{Voiceless sonorants}
The 2014 symbols $m̥$, $n̥$, $ŋ̊$, $l̥$ and $r̥$ are written according to their historical sources as a preinitial plus an ordinary sonorant.

\section{Spelling tests}
Representative forms and their Tibetan spellings appear in the web edition.

\end{document}
